We investigate whether a pretrained transformer's residual stream can adapt to a bijective token permutation---a systematic shuffling of which token IDs map to which surface forms---applied without any in-context demonstrations. Across experiments on \gpttwo and \pythia, we find that the model does \emph{not} recover ``modulo the shuffle.'' While cosine similarity between clean and ciphered hidden states increases at later layers, we show this is largely an architectural artifact: \layernorm drives all inputs---including random tokens---toward similar representations, with inter-text cosine similarity exceeding 0.96 in \gpttwo. A permutation-adjustment matrix designed to undo the cipher in embedding space makes similarity \emph{worse}, confirming the cipher's effect propagates non-linearly. The logit lens further shows the model never predicts the cipher-consistent next token at any layer, with mean ranks above 5{,}900 out of 50{,}257. \pythia, with weaker architectural convergence, reveals that ciphered representations sit 30\% above random but 30\% below clean at the final layer---partial structural extraction without genuine decoding. Our results demonstrate that the embedding lookup creates a persistent, non-invertible distortion: the forward pass cannot undo what amounts to a random permutation of 50{,}000 semantic embeddings.
